\subsection{A proved two-weight refinement using ambient homology}
\label{sec:two-weight-refinement}

The three-weight implementation uses only the boundary torus when choosing
its homology probes. The ambient manifold supplies additional information:
boundaries of proper surfaces occupy a one-dimensional subspace of the
torus homology. Exploiting that restriction gives the following proved
refinement. It is \emph{not implemented or benchmarked in this delivery};
the reported disc-query improvements use the implemented three-weight mode
and the stated comparison arms.

\begin{samepage}
\begin{proposition}[One ambient-certified boundary probe]
\label{prop:two-weight-disc}
Let $M$ be a compact connected orientable three-manifold with exactly one
torus boundary, and let $i:\partial M\hookrightarrow M$. There exists a
boundary cycle $\gamma$ with $i_*[\gamma]\ne0$ in $H_1(M;\F_2)$. For any
such cycle and every connected proper normal-surface component $C$,
\[
 C\text{ is a compressing disc}
 \quad\Longleftrightarrow\quad
 \chi(C)=1\ \text{ and }\ |\partial C\cap\gamma|\equiv1\pmod2.
\]
The intersection count can be computed by the same additive edge-point
weights as before. Thus two integer weights suffice after the indicated
ambient preprocessing.
\end{proposition}
\end{samepage}
\begin{proof}
All homology groups in this proof have coefficients in $\F_2$. Write
$b_j=\dim H_j(M)$. Poincar\'e--Lefschetz duality gives
$\chi(M,\partial M)=-\chi(M)$; combined with the Euler characteristic of
the pair, this gives $2\chi(M)=\chi(\partial M)=0$. Since $M$ is connected
with nonempty boundary, $b_0=1,b_3=0$, and hence $b_1-b_2=1$
\cite{Hatcher2002}.

In the long exact sequence of the pair, $H_0(\partial M)\to H_0(M)$ is an
isomorphism. Consequently $H_1(M)\to H_1(M,\partial M)$ is surjective.
Duality identifies the dimension of the latter group with $b_2$. Exactness
therefore gives $\dim\operatorname{im}i_*=b_1-b_2=1$. The kernel
\[
 L=\ker\bigl(H_1(\partial M)\longrightarrow H_1(M)\bigr)
\]
is a line in the two-dimensional torus homology. Every compact surface,
including a nonorientable one, has a relative fundamental class over
$\F_2$. Pushing that class into $(M,\partial M)$ shows that
$[\partial C]\in L$.

The torus intersection pairing is nondegenerate and alternating. A line
$L$ consequently has orthogonal complement $L$ itself. Since
$i_*[\gamma]\ne0$, the class $[\gamma]$ is outside $L$ and pairs to one
with its unique nonzero element. Odd intersection with $\gamma$ is therefore
equivalent to $[\partial C]\ne0$. The surface-classification and primitive
torus-slope argument already proved then gives the disc criterion. This
argument imposes neither irreducibility nor a torsion-free integral homology
assumption.
\end{proof}

There is a concrete finite preprocessing algorithm. Embed the two verified
boundary basis cycles into the ambient global edge-chain space over $\F_2$.
Form the span of boundaries of all ambient triangular face classes, counting
paired faces once and combining repeated edges modulo two. At least one
boundary basis cycle lies outside this span, because the boundary inclusion
has rank one. Choose such a cycle as $\gamma$. It is already a cycle, so
being outside the face-boundary span is exactly the required nonzero ambient
homology condition.

The preprocessing can carry a particularly small certificate: an ambient
cellular $1$-cocycle $u$ satisfying
\[
 u(\partial f)=0\quad\text{for every triangular face }f,
 \qquad u(\gamma)=1.
\]
These identities prove that $\gamma$ is not a boundary. A producer can find
$u$ by Gaussian elimination in $O(t^3)$ elementary operations over $\F_2$;
replay checks $O(t)$ constant-size face equations and the cycle evaluation,
without repeating elimination. The cycle and cocycle depend only on the
triangulation and can be reused across its candidate surfaces.

Both signatures have constant dimension. This refinement saves one
transported integer coordinate while adding ambient preprocessing; its practical benefit remains to be measured. It
does not contradict the earlier observation about two linear probes on the
\emph{unrestricted} torus class space. The new probe uses a certified
geometric restriction to the line $L$. With several boundary components the
rank argument above changes, so this particular one-probe conclusion must
not be transferred to that setting without a new analysis.
